\section{8. Analysis}

\subsection{8.1 Why Does ClusterFormer Work?}

\subsubsection{8.1.1 Semantic Grouping Reduces Redundancy}

\textbf{Figure 8.1} visualizes cluster assignments for a DistilBERT layer. 
Tokens with similar semantics (e.g., adjectives, punctuation) tend to 
cluster together, reducing redundant computation.

\subsubsection{8.1.2 Centroid-to-Key Preserves Global Context}

Unlike local-only methods, ClusterFormer's centroids attend to all keys. 
This captures long-range dependencies that windows would miss.

\subsubsection{8.1.3 Local Residual Fills the Gap}

The sliding-window branch handles fine-grained syntactic structure. The 
learned blend $\alpha$ adapts per layer.

\subsection{8.2 Complexity in Practice}

\textbf{Figure 8.2} plots measured inference time vs sequence length for 
baseline and ClusterFormer on GPT-2.

\begin{table}[H]
\centering
\begin{tabular}{llll}
\toprule
\textbf{$L$} & \textbf{Baseline (ms)} & \textbf{ClusterFormer (ms)} & \textbf{Speedup} \\
\midrule
64 & 2.1 & 2.4 & 0.88× \\
128 & 4.2 & 3.8 & 1.10× \\
256 & 12.5 & 6.9 & 1.81× \\
512 & 44.2 & 15.3 & **2.89×** \\
1024 & 172.8 & 32.1 & **5.38×** \\
\bottomrule
\end{tabular}
\end{table}

\textbf{Crossover point} is around $L \approx 100$: below that, ClusterFormer 
overhead exceeds savings.

\subsection{8.3 Cluster Quality Analysis}

\textbf{Figure 8.3} shows the distribution of tokens per cluster. Well-balanced 
clustering has near-uniform distribution.

\textbf{Table 8.1: Cluster balance metrics.}

\begin{table}[H]
\centering
\begin{tabular}{llll}
\toprule
\textbf{Model} & \textbf{Min tokens/cluster} & \textbf{Max tokens/cluster} & \textbf{Ideal} \\
\midrule
DistilBERT & 5 & 12 & 8 \\
RoBERTa & 4 & 13 & 8 \\
GPT-2 & 5 & 12 & 8 \\
\bottomrule
\end{tabular}
\end{table}

\subsection{8.4 Comparison with Performer}

\textbf{Table 8.2: ClusterFormer vs Performer (theoretical).}

\begin{table}[H]
\centering
\begin{tabular}{lll}
\toprule
\textbf{Feature} & \textbf{Performer} & \textbf{ClusterFormer} \\
\midrule
Complexity & $\mathcal{O}(L m d)$ & $\mathcal{O}(L K d)$ \\
Approximation & Random features & Exact on centroids \\
Fidelity & Depends on $m$ & Depends on $K$ \\
Kernel assumption & None (universal) & Clustering assumption \\
Local context & None & Yes (window) \\
\bottomrule
\end{tabular}
\end{table}

\subsection{8.5 Ablation Insights}

\textbf{Table 8.3: Cumulative component contributions.}

\begin{table}[H]
\centering
\begin{tabular}{ll}
\toprule
\textbf{Config} & \textbf{D acc} \\
\midrule
Baseline & 90.71\% \\
+ Frozen base + LoRA & 88.5 \\
+ Centroid attention (no local) & 88.8 \\
+ Local window & 89.7 \\
+ Learnable temp/blend & 89.91 \\
\bottomrule
\end{tabular}
\end{table}

\subsection{8.6 Numerical Stability}

\subsubsection{8.6.1 Empty Clusters}

Handled by masking logits to $-\infty$ before softmax.

\subsubsection{8.6.2 NaN in cdist}

\texttt{torch.nan_to_num} guards against numerical issues.

\subsubsection{8.6.3 Temperature Clamping}

$\tau$ is clamped to $[0.1, 10.0]$ to avoid degenerate softmax.

\subsection{8.7 Failure Analysis}

No systematic failures observed.

\subsection{8.8 Summary of Insights}

1. \textbf{Semantic clustering works} — tokens group meaningfully.
2. \textbf{Global + local is essential} — either alone underperforms.
3. \textbf{Crossover at $L \approx 100$} — beneficial for long sequences.
4. \textbf{Learned $\tau$ and $\alpha$ help} — small gains from learnable scalars.
5. \textbf{Not a compression method} — parameter count increases.